% =====================================================================
% Appendix: the full audit table -- SKELETON.
% The appendix table is generated from results/audit_verdicts_rev2.csv; do not
% hand-transcribe it. Column semantics below are copied from that file's header
% and from docs/AUDIT_VERDICTS_2026-07-20.md.
% =====================================================================

\section{Full audit table}
\label{app:audit-table}

\TODO{typeset the full 17-row table directly from
\texttt{results/audit\_verdicts\_rev2.csv} (one row per claim, landscape, small
font). Do not retype values.}

The released CSV carries, per claim: identity and frame; method family, bit
width and benchmark; the reported or imputed $n$ with its \texttt{n\_basis};
the claimed margin with its \texttt{margin\_basis}; the imputed discordance with
its match tier and the number of atlas cells behind it; indeterminacy status,
kind, and reason; the components that remain computable for indeterminate
claims; the V3 reproducibility value; V1 as paired and independent-binomial MDD
with both margin ratios; V2 required-$n$ at 1, 2 and 3\,pp under both variance
models, plus at the claim's own margin; the applicable margin and its basis;
the margin-sensitivity flag; and both verdicts
(\texttt{verdict} and \texttt{verdict\_at\_registered\_2pp}).

\paragraph{Reading the transparency columns.} Rows marked indeterminate carry
numbers in some V1/V2 columns. Those numbers are \emph{not} verdicts and are not
counted in $K$; they are retained so that a reader can see exactly what the
analysis could and could not evaluate. R04's GSM8K computation is the largest
such quantity and is discussed in \S\ref{sec:audit:r04}.

\section{Audit method detail}
\label{app:audit-method}

% RELOCATED 2026-07-27 (phase 4 trim, candidates 1 and 3). Moved VERBATIM in
% content from sections/audit.tex; nothing was deleted in the move. The body
% keeps both findings and points here for the mechanism.

\subsection{Discordance imputation}
\label{app:audit:imputation}

% SOURCE: docs/AUDIT_VERDICTS_2026-07-20.md §Method, imputation-tier table.
The discordance rate $p_d$ is not reported by any audited source, so it is
imputed from the atlas (\S\ref{sec:atlas}) by matching the nearest (method
family, bit width, benchmark) cell, most-specific tier first, taking the
\emph{median} over the first non-empty tier because per-cell discordance is
right-skewed. One claim matched at tier~1 (family + bits + benchmark), 11 at
tier~2 (family + bits), 2 at tier~3 (bits + benchmark), 1 at tier~4 (bits), and
2 fell through to the global tier. A tier whose target field is \texttt{None}
cannot match, so a pruning claim with no bit width descends automatically rather
than being forced into a wrong cell. Both disclosed feasibility-probe pairs (99
atlas cells) are excluded from the imputation pool, per the atlas registration.

\subsection{The two robustness readings, and why they stay separate}
\label{app:audit:robustness}

% SOURCE: docs/AUDIT_VERDICTS_2026-07-20.md §Headline (secondary reading) and
% §"Margin-sensitive (1 of 12 determinate)"; results/audit_verdicts_rev2.csv
% columns verdict_at_registered_2pp and margin_sensitive.
% CARE -- THE TWO "1 of 12" FIGURES ARE DIFFERENT FACTS. First: one determinate
% claim is underpowered under the uniform 2 pp yardstick. Second: one
% determinate claim's verdict is unstable across the 1 pp -> 3 pp sweep. They
% happen to concern the same claim (R01) and MUST NEVER be merged into a single
% sentence or read as one finding. The body states both separately and points
% here; this appendix carries the reasoning.

Both readings ship in the released CSV (\texttt{verdict},
\texttt{verdict\_at\_registered\_2pp}). The factor-of-four difference between
judging each claim against its own stated margin and judging all of them against
the uniform registered 2\,pp margin is \emph{interpretive, not statistical}, and
the reasoning that selects the primary reading is at
\S\ref{sec:prereg:choices}: the frozen label is ``underpowered for its
\emph{own} assertion'', and judging a 0.15\,pp parity claim against a 2\,pp
yardstick audits a claim nobody made.

Margin sensitivity is a separate property, and it qualifies a headline verdict,
so it is counted over the claims that have one. R01's verdict flips between
underpowered and adequately powered across the 1\,pp\,$\to$\,3\,pp sweep, making
that verdict an artefact of margin choice rather than a robust finding; it is
reported as such rather than counted as evidence in either direction. R04 and
R14 also flip across the sweep and retain the \texttt{margin\_sensitive} column
in the released CSV, but neither carries a verdict to qualify.

These are two distinct properties---a verdict under an alternative yardstick,
and instability of a verdict across the sweep---that coincide on one claim.

\section{Certification tables at 1\,pp and 3\,pp}
\label{app:certification-margins}

\TODO{typeset the 1\,pp and 3\,pp margins from
\texttt{results/certification\_tables\_rev2.csv}; the 2\,pp margin is
Table~\ref{tab:certification} in the main text. Note the quadratic margin
scaling: MMLU's median requirement is 8{,}656 items at 1\,pp, 2{,}164 at 2\,pp,
and 962 at 3\,pp.}
% SOURCE: results/certification_tables_rev2.csv, rows mmlu at margin_pp
% 1.0/2.0/3.0, column required_n_median.
%
% CORRECTION (2026-07-22). This block previously read "944 at 3 pp" and pointed
% at the rev-1 CSV, while its 1 pp and 2 pp values (8,656 and 2,164) were
% already rev-2. The triple was mixed: two of three values had been updated and
% the third had not. Commit 272136b lists the intended change as "944 -> 962",
% so the value was identified at the time and missed in this one place.
% Corrected to 962 (results/certification_tables_rev2.csv, mmlu, margin_pp 3.0,
% required_n_median), and the SOURCE pointer moved to the rev-2 file to match.
% The stale value sat inside this TODO rather than in typeset body text, so it
% was never a rendered claim -- but it is the instruction this appendix would
% have been typeset from.

\subsection{Scope caveats carried from the main text}
\label{app:cert:caveats-detail}

% RELOCATED 2026-07-27 (phase 4 trim, candidate 6) from
% sections/certification.tex §"Scope and caveats", caveats 2, 3 and 4, moved
% VERBATIM in content. Caveats 1 and 5 remain in the body.
% SOURCE: docs/CERTIFICATION_TABLES_2026-07-20.md §"Scope and caveats";
% results/certification_tables_rev2.csv column n_atlas_cells.

\paragraph{Cell counts behind each family.} Rev-2 moved
\texttt{arc\_challenge} to 17 cells, \texttt{gsm8k} to 24, \texttt{hellaswag}
to 23 and \texttt{winogrande} to 23, taking those out of the thin category;
\texttt{mmlu} (1{,}311), \texttt{bbh} (192) and \texttt{math} (56) are well
supported. Only \texttt{mmlu\_pro} (5) and \texttt{ifeval} (8) remain thin, as
the main text states.

\paragraph{Family aggregation mixes two sources of variation.} The atlas
collapses MMLU's 57 per-subject cells and BBH/MATH/MuSR/GPQA's per-subtask cells
into families, so a family's spread combines subject-level with model-level
variation. This widens the p25--p75 band relative to what a single practitioner
evaluating a single model would see, and therefore makes the quartile columns
\emph{conservative} rather than optimistic.

\paragraph{Feasibility-probe pairs are excluded.} Both disclosed
feasibility-probe pairs (99 cells) are excluded per the atlas registration §6,
as tiny hand-built sanity pairs---$n$ as low as 10, discordance up to
0.9---that would distort every quartile.

\section{Preregistration documents}
\label{app:registrations}

\TODO{reproduce, verbatim, \texttt{PREREGISTRATION.md} and the three
2026-07-15 registrations with their dated amendments, or link them at the
archived commit. Reviewers should be able to read the frozen rule text that
\S\ref{sec:prereg} describes.}
